\chapter{Revisão Bibliográfica}
\label{chap:revisao}

Esta revisão está organizada em quatro eixos, cada um ligado a uma decisão
metodológica adotada no relatório: a escolha da técnica de verificação, a
forma de representar a rede neural quantizada, a origem do controlador
avaliado e o tratamento dado a componentes de Inteligência Artificial
Generativa.

\section{Bounded Model Checking e verificação de software em C}
O \textit{Bounded Model Checking} (BMC) verifica uma propriedade
explorando as execuções de um programa até um número finito de passos,
codificando o programa e a negação da propriedade como uma fórmula e
delegando a busca por uma violação a um \textit{solver}. Biere et al.
\cite{biere1999symbolic} propuseram o BMC como alternativa ao
\textit{model checking} simbólico baseado em diagramas de decisão binários,
mostrando que a codificação proposicional de um número limitado de passos
já é suficiente para encontrar contraexemplos em muitos casos práticos.
Cordeiro, Fischer e Marques-Silva \cite{cordeiro2012smt} estenderam essa
ideia para software ANSI-C incorporado, propondo a verificação baseada em
SMT (\textit{Satisfiability Modulo Theories}) em vez de lógica proposicional
pura, o que permite raciocinar diretamente sobre aritmética de ponteiros,
inteiros de largura fixa e ponto flutuante -- exatamente o tipo de código
produzido pelos harnesses C usados neste relatório. O ESBMC evoluiu dessa
linha de trabalho; Gadelha, Ismail e Cordeiro \cite{gadelha2018esbmc}
descrevem a ferramenta em sua versão 5.0 como um verificador de força
industrial que combina múltiplos \textit{solvers} SMT, oferece suporte a
concorrência e mantém compatibilidade com o padrão C. Para este projeto,
a relevância dessa linha está em três decisões diretas: usar C como
linguagem intermediária entre o modelo treinado e o verificador, expressar
propriedades como pares \texttt{assume}/\texttt{assert} e aceitar que
\texttt{TIMEOUT} é um resultado epistemicamente distinto de prova ou
refutação, e não um indício de segurança.

\section{Verificação formal de redes neurais e redes quantizadas}
Huang et al. \cite{huang2017robustness} formularam a verificação de
robustez de redes neurais profundas como uma busca sistemática por
perturbações de entrada que alterem a classificação, discutindo garantias
de cobertura da vizinhança de uma entrada. Esse trabalho estabeleceu que
verificar uma rede neural exige tratar explicitamente a rede como um
programa com estrutura conhecida (camadas, pesos e ativações), e não como
uma caixa-preta amostrada por testes. Gopinath et al.
\cite{gopinath2021verifying} tornaram esse problema mais próximo do
presente relatório ao tratar redes neurais \emph{quantizadas}: os autores
mostram que a verificação baseada em SMT precisa modelar explicitamente a
aritmética de precisão finita -- escala, arredondamento e saturação --
porque uma propriedade provada para a rede em ponto flutuante não se
transfere automaticamente para sua versão quantizada. Meng et al.
\cite{meng2022qnnverifier} consolidaram esse fluxo na ferramenta
QNNVerifier, que converte um modelo treinado em um programa C anotado com
propriedades e o submete a um \textit{model checker} baseado em SMT. O
harness Q8.8 e o desenho da comparação Float32~$\times$~Q8.8 deste
relatório seguem diretamente essa mesma lógica: o controlador quantizado é
reescrito como código C, e a fidelidade em relação à versão original é
tratada como uma propriedade a ser medida e verificada, não como uma
suposição.

\section{Controle por aprendizado por reforço e a planta CartPole}
O controlador analisado neste relatório é treinado com \textit{Deep
Deterministic Policy Gradient} (DDPG), algoritmo de aprendizado por reforço
com espaço de ações contínuo proposto por Lillicrap et al.
\cite{lillicrap2016ddpg}, que combina uma rede de ator (a política, cuja
representação quantizada é o objeto verificado neste relatório) com uma
rede de crítico usada apenas durante o treinamento. A planta de controle
usada como estudo de caso, o \textit{Cart-Pole}, é um problema clássico de
controle não linear com realimentação: Barto, Sutton e Anderson
\cite{barto1983cartpole} formalizaram o problema de equilibrar um pêndulo
invertido sobre um carrinho por meio de elementos adaptativos, e essa
formulação de estado (posição e velocidade do carrinho, ângulo e
velocidade angular do pêndulo) é a mesma usada para definir o domínio
operacional e as propriedades de segurança verificadas neste relatório.

\section{Componentes de Inteligência Artificial Generativa}
O título do projeto também contempla sistemas de Inteligência Artificial
Generativa. Diferentemente dos três eixos anteriores, não há neste
relatório um estudo de caso de IA generativa com artefato executável,
comando registrado e resultado auditável equivalente ao do controlador
CartPole; por isso, esta revisão não apresenta uma linha de trabalhos
específicos para esse eixo. A posição adotada é metodológica: os mesmos
princípios de verificação -- expressar a propriedade como
\texttt{assume}/\texttt{assert} sobre uma representação em C, delegar a
busca ao ESBMC e distinguir prova, refutação e resultado inconclusivo --
são, em princípio, aplicáveis a código produzido ou auxiliado por modelos
de linguagem, mas essa transferência é tratada neste relatório como trabalho futuro
condicionado à existência de um caso executável, comando registrado e
saída verificável (PQ5, discutida nos Capítulos~\ref{chap:objetivos} e
\ref{chap:resultados}).
